% ============================================================
\section{Ambiente di Lavoro}
% ============================================================

\subsection{Repository e Versione del Progetto}

Il progetto Apache OpenJPA è stato ottenuto tramite fork del repository ufficiale su GitHub. Il lavoro si è svolto sul branch \texttt{release-4.1.1}, corrispondente al tag di rilascio \texttt{4.1.1}. L'ambiente di sviluppo primario è Windows 11 con WSL2 (Ubuntu); tutti gli script e le configurazioni sono stati mantenuti compatibili con ambienti Linux e macOS.

\subsection{Versione Java}

OpenJPA 4.1.1 è progettato per Java 17. L'utilizzo di Java 26 causa errori di compilazione nel modulo \texttt{openbooks} a causa di modifiche alle regole di type inference dei generics introdotte nelle versioni più recenti del compilatore Java. Per questa ragione, Java 17 è stato adottato come versione principale per tutte le operazioni di build e testing. Lo switch di versione viene effettuato tramite variabile d'ambiente:

\begin{verbatim}
export JAVA_HOME=/usr/lib/jvm/java-17-openjdk-amd64
export PATH=$JAVA_HOME/bin:$PATH
\end{verbatim}

Lo strumento EvoSuite, utilizzato per la generazione automatica di test, richiede anch'esso Java 17, risultando pertanto compatibile con questa configurazione.

\subsection{Rimozione dei Test Originali}

Al fine di sostituire interamente la suite di test con quella sviluppata nel presente progetto, i sorgenti di test originali sono stati rimossi da tutte le directory \texttt{src/test} dei sottomoduli, tramite il comando:

\begin{verbatim}
find <project-root> -type d -path "*/src/test" \
    -exec sh -c 'rm -rf "$1"/*' _ {} \;
\end{verbatim}

\subsection{Adeguamento dei File \texttt{pom.xml}}

La rimozione dei contenuti di \texttt{src/test} ha reso non valide alcune \texttt{<execution>} Maven definite nei \texttt{pom.xml} di diversi sottomoduli. Tali execution erano legate all'enhancement bytecode delle entità JPA di test e referenziavano file (quali \texttt{persistence.xml} ed \texttt{enhancer.xml}) presenti in \texttt{src/test}. Le fasi coinvolte sono \texttt{process-test-classes} e \texttt{test-compile}, che precedono l'esecuzione di Surefire e non sono controllabili tramite il flag \texttt{-DskipTests}.

Sono stati identificati quattro tipi distinti di execution da disabilitare, ciascuno con una causa e un impatto funzionale diverso:

\begin{itemize}
    \item \textbf{Enhancement tramite Ant} (\texttt{maven-antrun-plugin}) — diversi moduli eseguivano script Ant (tipicamente \texttt{enhancer.xml}) in fase \texttt{process-test-classes} per applicare il bytecode enhancement alle entity JPA di test. Questi script cercavano il file \texttt{persistence.xml} in \texttt{src/test/resources}; la sua assenza causava un errore \textit{MetaDataFactory could not be configured}. Commentando le execution corrispondenti si rimuove l'enhancement delle classi di test originali, non rilevante per la campagna di testing sviluppata.

    \item \textbf{Enhancement tramite plugin nativo} (\texttt{openjpa-maven-plugin}, goal \texttt{test-enhance}) — alcuni moduli utilizzavano il plugin Maven nativo di OpenJPA in sostituzione di Ant per lo stesso scopo. Il comportamento e l'errore prodotto sono identici a quelli del caso precedente; la soluzione adottata è analoga.

    \item \textbf{Spacchettamento di classi di supporto ai test} (\texttt{maven-dependency-plugin}, goal \texttt{unpack}) — in fase \texttt{test-compile}, alcuni moduli estraevano da un test-jar di \texttt{openjpa-persistence-jdbc} un insieme di classi di utilità per i test originali (package \texttt{org.apache.openjpa.persistence.test}). Con la rimozione dei test originali, queste classi non sono più necessarie e la loro estrazione è stata disabilitata.

    \item \textbf{Build di sotto-progetti di integrazione} (\texttt{maven-invoker-plugin}) — il modulo \texttt{openjpa-tools/openjpa-maven-plugin} eseguiva build Maven separati su sotto-progetti in \texttt{src/it} per verificare il corretto funzionamento del plugin Maven di OpenJPA come strumento esterno. Questi sotto-progetti dipendevano da risorse di test non più presenti. Il loro scopo è testare il plugin stesso come tool di build, non le entity del progetto; la loro disabilitazione non influisce sulla capacità di eseguire test unitari o di integrazione tramite Surefire e Failsafe.
\end{itemize}

In totale, le execution commentate interessano dieci sottomoduli distribuiti tra il kernel, i moduli di integrazione e gli esempi del progetto. In tutti i casi, l'intervento è limitato al blocco \texttt{<execution>} del \texttt{pom.xml} del modulo interessato, lasciando invariata la configurazione dei plugin Surefire e Failsafe responsabili dell'esecuzione dei test.
